\chapter{Couche 0 — Le socle}

Avant la première ligne de code métier, trois outils décident de ce que le
projet pourra devenir. On les oublie vite parce qu'ils ne « font » rien de
visible, et ce sont pourtant eux qu'on regrette quand ils manquent.

% -----------------------------------------------------------------------------
\outil{Git}{brique 1 · versionnement}

\quoi
Un système qui enregistre l'historique complet d'un dossier : qui a changé
quoi, quand, et pourquoi. Chaque enregistrement --- un \emph{commit} --- est
une photographie du projet entier, pas seulement des fichiers modifiés.

\pourquoi
L'alternative honnête, c'est de copier le dossier en
\cmd{projet-v2-final-vraiment}. Ça marche jusqu'au jour où l'on veut savoir
quand un bug est apparu, ou revenir trois jours en arrière sans perdre le
travail d'hier.

Il existe d'autres systèmes --- Subversion, Mercurial --- mais Git est celui que
toute l'industrie utilise, et surtout le seul que les plateformes
d'intégration continue savent lire nativement. Choisir autre chose aurait coupé
le projet de GitHub Actions, donc de toute la couche 6.

\role
Trente-cinq commits, chacun avec un message qui explique \emph{pourquoi} et pas
seulement \emph{quoi}. Les branches servent à isoler un travail en cours : la
branche \cmd{aziz-youssef-modif} porte l'intégration d'une contribution
extérieure sans toucher à \cmd{main} tant qu'elle n'est pas vérifiée.

Un fichier \fichier{.gitignore} tient hors du dépôt ce qui ne doit pas y
entrer : les clés (\cmd{.secret\_key}, \cmd{.vault\_key}), le fichier
\cmd{.env} des secrets de la pile, les environnements virtuels, et les archives
de travail des utilisateurs.

\apport
Sans Git, aucune des briques de la couche 6 n'existe : pas d'intégration
continue, pas de recherche de secrets dans l'historique, pas de déploiement
déclenché par un envoi de code. C'est la brique dont toutes les autres
dépendent indirectement.

% -----------------------------------------------------------------------------
\outil{Python 3.10}{brique 2 · langage}

\quoi
Le langage dans lequel est écrite toute la logique du projet : l'application
web, le démon de nettoyage, l'agent des nœuds, les scripts de génération de
schémas.

\pourquoi
Trois contraintes ont guidé le choix.

D'abord la \textbf{bibliothèque standard}. L'agent qui tourne sur chaque worker
ne doit rien avoir à installer : il utilise \cmd{urllib} pour les appels HTTP,
\cmd{json}, \cmd{ssl}, \cmd{subprocess}, et lit \cmd{/proc/meminfo}
directement. En Node.js ou en Go, il aurait fallu déployer un binaire ou un
environnement sur chaque nœud.

Ensuite l'\textbf{écosystème} : Flask, SQLAlchemy, Ansible lui-même sont en
Python. Rester dans un seul langage évite d'avoir deux chaînes d'outillage.

Enfin la \textbf{version} : 3.10 précisément, parce que la cible de
déploiement est Ubuntu 22.04, qui fournit cette version. Ce détail nous a
coûté une panne : nous avions épinglé une bibliothèque qui n'existait pas pour
3.10, les tests passaient sur nos postes et le déploiement échouait.

\begin{piege}{La version du poste n'est pas celle de la cible}
C'est l'erreur la plus banale et la plus coûteuse. La chaîne d'intégration
force désormais \cmd{PYTHON\_VERSION: 3.10} : elle teste avec le Python de la
machine de production, pas celui du développeur.
\end{piege}

\role
Environ 3\,500 lignes pour l'application, 1\,600 pour les tests, plus les
scripts d'outillage. Le code est séparé en modules dont chacun ignore les
couches au-dessus de lui : \fichier{docker\_ops.py} ne sait pas que Flask
existe, ce qui le rend testable sans serveur.

\apport
Un seul langage du navigateur au worker, et un agent déployable sans rien
installer.

% -----------------------------------------------------------------------------
\outil{venv et requirements épinglés}{brique 3 · reproductibilité}

\quoi
Un \emph{environnement virtuel} est un dossier contenant une copie isolée de
Python et de ses bibliothèques. Le fichier \fichier{requirements.txt} liste les
bibliothèques avec leur \textbf{version exacte} : \cmd{Flask==3.1.3}, pas
\cmd{Flask}.

\pourquoi
Sans isolation, installer une bibliothèque pour ce projet peut casser un autre
projet sur la même machine. Sans versions figées, le projet marche aujourd'hui
et casse demain parce qu'une dépendance a publié une version incompatible.

L'alternative moderne serait Poetry ou uv, plus confortables. On est resté sur
\cmd{venv} et \cmd{pip}, qui sont dans la bibliothèque standard : aucune
installation préalable n'est nécessaire sur la machine cible, et le rôle Ansible
peut créer l'environnement avec ce que fournit Ubuntu.

\role
Deux fichiers : \fichier{requirements.txt} pour ce qui tourne en production,
\fichier{requirements-dev.txt} pour l'outillage de tests et d'analyse. Le rôle
Ansible \cmd{webapp} installe exactement le premier sur le contrôleur --- pas
une liste recopiée à la main, ce qui a déjà été une source de divergence.

\apport
C'est ce qui permet à \cmd{pip-audit} d'exister (couche 6) : on ne peut
confronter des dépendances à une base de vulnérabilités que si l'on sait
exactement lesquelles on utilise.

% -----------------------------------------------------------------------------
\section*{Ce que la plateforme sait faire après la couche 0}

Rien d'utilisable. Mais le projet est désormais \emph{reproductible} : on peut
le cloner, créer l'environnement, et obtenir exactement les mêmes versions. Sans
cela, toutes les couches suivantes reposeraient sur du sable.
